\section{Case study: baseline prototype-target scenario}
\label{sec:case-study}

The preceding model results describe the individual thermal, hydraulic,
pressure-integrity, and potting-process relationships within their stated
limits. This case study combines them with the tube-supply-cost criterion using
one internally consistent set of example values. The values in
Table~\ref{tab:design_boundary_assumptions} define this example; they are not
design allowables or universal feasibility limits.

\subsection{Example design-boundary map}
\label{sec:case-design-boundary}

Figure~\ref{fig:design_boundary} combines the preceding thermal, hydraulic,
mechanical, process, and tube-supply-cost constraints in one example design-space
view. The baseline values and thresholds are summarized in
Table~\ref{tab:design_boundary_assumptions}.

\begin{table}[htbp]
\caption{Baseline thresholds used for the design-boundary map.}
\label{tab:design_boundary_assumptions}
\centering
\small
\setlength{\tabcolsep}{3pt}
\renewcommand{\arraystretch}{0.95}
\begin{tabular*}{\textwidth}{@{}ll@{\extracolsep{\fill}}ccc@{}}
\toprule
Boundary & Scope & Symbol & Limit & Unit \\
\midrule
Min. wall thickness & Al, drawn tube
  & $t$ & $\ge 0.070$ & $\mathrm{mm}$ \\
Min. wall thickness & PA, precision route
  & $t$ & $\ge 0.025$ & $\mathrm{mm}$ \\
Coolant throughput & fixed velocity model
  & $\dot{V}_{\mathrm{c}}$ & $\ge 10$ & $\mathrm{L\,min^{-1}}$ \\
Pressure drop & fixed friction model
  & $\Delta p_{\mathrm{i},\mathrm{fric}}$ & $\le 0.5$ & $\mathrm{bar}$ \\
Tube-supply cost & cost index
  & $C_{\mathrm{tube}}/C_{\mathrm{fin}}$ & $<5$ & -- \\
Burst pressure & normal extrusion
  & $p_{\mathrm{b},\mathrm{tol}}$ & $\ge 6$ & $\mathrm{bar}$ \\
Capillary rise & potting, $G=10$
  & $h_{\mathrm{cap}}$ & $\le 2$ & $\mathrm{mm}$ \\
\bottomrule
\end{tabular*}
\end{table}
\FloatBarrier

Under this example value set, cost, pressure-integrity, potting, and
coolant-side constraints remove large parts of the thermally attractive design
space. The remaining PA region defines a prototype target window for this value
set; turning it into a design rule would require data over production
tolerances, potting quality, and pressure cycling.
Under the example normal-route tolerance, the PA minimum-wall line itself
fails the $6\,\mathrm{bar}$ burst criterion:
$t_{\mathrm{min},\mathrm{PA}}=0.025\,\mathrm{mm}$ leaves a minimum local wall of
$t_{\mathrm{loc},\mathrm{min}}=5\,\mu\mathrm{m}$.
The PA region accepted by the combined constraints requires walls moderately
above the example process minimum,
$t \gtrsim \Delta t + 0.03\,d_{\mathrm{o}}$, about
$0.030$--$0.040\,\mathrm{mm}$ at sub-millimeter diameters.
Its lower diameter edge in this scenario is not set by the capillary limit.
Under the example $G=10$ benchmark-pitch convention and the
$h_{\mathrm{cap}}\le2\,\mathrm{mm}$ limit, the updated PA/process capillary
limit reaches its threshold near $d_{\mathrm{o}}\approx 0.27\,\mathrm{mm}$,
below the tube-supply-cost cutoff near $d_{\mathrm{o}}\approx 0.35\,\mathrm{mm}$
and the straight-tube friction limit near $d_{\mathrm{o}}\approx
0.31\,\mathrm{mm}$.
For aluminum, the case-study minimum-wall boundary starts at
$d_{\mathrm{o}}\approx 0.63$--$0.65\,\mathrm{mm}$ because the example
tube-supply-cost cutoff excludes smaller aluminum diameters.

\begin{figure}[htbp]
\centering
\includegraphics[width=\textwidth]{20_design_boundary_lines_alu_pa}
\caption{Design-boundary map for Al and PA tubes.
  Color encodes $k_{\mathrm{o}}A_{\mathrm{o}}$ for constraint-admitted points;
  boundary colors follow the legend and hatching marks excluded sides. Blue
  cross: validated aluminum benchmark. The $d_{\mathrm{o}}$ axis is logarithmic.}
\label{fig:design_boundary}
\end{figure}

\FloatBarrier
\subsection{Feasible PA-to-aluminum conductance comparison}
\label{sec:case-delta}

The comparison in Fig.~\ref{fig:tech_adj} uses the combined feasibility masks
from Table~\ref{tab:design_boundary_assumptions}. For
$m\in\{\mathrm{Al},\mathrm{PA}\}$, the mask
$\mathcal{F}_m(d_{\mathrm{o}},\tau)\in\{0,1\}$ is set to 1 when the design
point satisfies all applicable thresholds in
Table~\ref{tab:design_boundary_assumptions}, and to 0 otherwise.
For each feasible PA point, the aluminum comparator is the nearest feasible
aluminum geometry at the same wall-thickness ratio $\tau$ and at an equal or
larger outer diameter. 

\begin{figure}[htbp]
\centering
\includegraphics[width=\textwidth]{22_tech_adjusted_kA_change}
\caption{Case-study relative conductance difference
  $\Delta (k_{\mathrm{o}}A_{\mathrm{o}})^{\,\mathrm{feas}}$.
  Color marks PA points passing all example constraints.
  Each PA point is compared with the nearest constraint-admitted aluminum diameter at
  the same $\tau$; uncolored areas are excluded or have no admitted aluminum
  comparator. Solid black curves are isolines of
  $\Delta (k_{\mathrm{o}}A_{\mathrm{o}})^{\,\mathrm{feas}}$.
  Dashed dark-blue and dotted dark-green curves mark the Al and PA combined-constraint
  lower-diameter boundaries, respectively. Blue cross: validated aluminum
  benchmark geometry. The $d_{\mathrm{o}}$ axis is logarithmic.}
\label{fig:tech_adj}
\end{figure}

The feasible PA-to-aluminum comparison is defined by
Eq.~\eqref{eq:feasible_pa_al_comparison},
\be
  \begin{aligned}
  d_{\mathrm{o},\mathrm{Al}}^{*}(d_{\mathrm{o}},\tau)
  &=
  \min\left\{
    \tilde d_{\mathrm{o}}\ge d_{\mathrm{o}}\; \middle|\;
    \mathcal{F}_{\mathrm{Al}}(\tilde d_{\mathrm{o}},\tau)=1
  \right\},\\
  \Delta (k_{\mathrm{o}}A_{\mathrm{o}})^{\,\mathrm{feas}}
  &=
  \frac{
    (k_{\mathrm{o}}A_{\mathrm{o}})_{\mathrm{PA}}(d_{\mathrm{o}},\tau)
  }{
    (k_{\mathrm{o}}A_{\mathrm{o}})_{\mathrm{Al}}
    (d_{\mathrm{o},\mathrm{Al}}^{*},\tau)
  }
  - 1,
  \qquad
  \mathcal{F}_{\mathrm{PA}}(d_{\mathrm{o}},\tau)=1 .
  \end{aligned}
  \label{eq:feasible_pa_al_comparison}
\ee

In this map a positive value marks a PA point that exceeds the nearest
constraint-admitted aluminum option at the same wall-thickness ratio, whereas a
negative value marks a point where aluminum remains superior. The relevant PA
target is the conditional sub-millimeter interval identified in
Sect.~\ref{sec:case-design-boundary}, where no constraint-admitted aluminum
comparator exists at a smaller diameter. The comparison isolates
geometry-dependent conductance; burst/leak tests and bundle heat-rate and
pressure-drop tests would have to show prototype qualification and system
performance.

\FloatBarrier
